% Propietario conservado: /Users/ruben/Documents/ChatGPT/jueces y controles/REVISION_ARGUMENTAL_APERTURAS_CIERRES_20260915/II/manuscript_es/sections/10b_boltzmann.tex
\section{Constante de Boltzmann, information et transduction thermique}
\label{ii:sec:ii-boltzmann}

\subsection{Structure dimensionnelle de la transduction}

La définition structurelle de la constante de Boltzmann utilise les
droites de température et d'énergie. La transduction positive
\[
 k_B:\mathcal L_\Theta\longrightarrow\mathcal L_E
\]
associe à une température son échelle énergétique. Dans la normalisation
tritique, la section de température du cycle satisfait
\[
 k_B\Theta_{\rm clk}=\frac{E_P}{\log3}.
\]
Le membre de droite est fixé par l'action et la période préalablement
construites. Le facteur \(\log3\) procède du contenu informationnel
d'une distribution uniforme sur les trois états locaux.

Soient \(u_E,u_\Theta\) des bases positives et écrivons
\(E_P=\epsilon_Pu_E\), \(\Theta_{\rm clk}=\theta_{\rm clk}u_\Theta\)
et \(k_B(u_\Theta)=b\,u_E\). Alors
\[
 b\,\theta_{\rm clk}=\frac{\epsilon_P}{\log3}.
\]
Pour \(u'_E=s_Eu_E\) et \(u'_\Theta=s_\Theta u_\Theta\),
les coordonnées se transforment selon
\[
 b'=\frac{s_\Theta}{s_E}b,\qquad
 \theta'_{\rm clk}=\frac{\theta_{\rm clk}}{s_\Theta},\qquad
 \epsilon'_P=\frac{\epsilon_P}{s_E}.
\]
L'identité énergétique demeure covariante. Cette formulation réunit
la structure individuelle du transducteur et la normalisation de la
section sur laquelle il agit.

\begin{remark}[Coordonnée thermique et choix des bases]
Le transducteur et l’identité énergétique précédente sont covariants.
Une valeur numérique individuelle de \(k_B\) se rapporte aux bases positives
\(u_E,u_\Theta\) ; sa transformation est \(b'=(s_\Theta/s_E)b\).
Par conséquent, la sélection interne d’une coordonnée numérique nécessite,
en outre, une construction qui fixe ces bases. L’identité énergétique
ne sélectionne pas à elle seule un tel choix. La coordonnée métrologique
ne définit ni l’action ni la période.
\end{remark}

La distinction aide à interpréter l'entropie. La fonctionnelle
\(s(\rho)=-\operatorname{Tr}(\rho\log\rho)\) est adimensionnelle.
L'entropie dimensionnelle est \(S(\rho)=k_Bs(\rho)\), comprise dans
la droite correspondante. La quantité \(\Theta S\) appartient à
la droite d'énergie. On conserve ainsi le sens mathématique des
égalités qui réunissent information, température et action.

En particulier, l'information de continuation
\eqref{ii:eq:ii-entropia-continuaciones} admet la publication dimensionnelle
\(S_n^{\rm gen}(x)=k_Bs_n(x)\). Les branches admissibles et leur mesure
déterminent \(s_n\) ; le transducteur fixe l'échelle d'entropie. Pour
une distribution uniforme sur trois possibilités, la valeur est
\(k_B\log3\). Lorsque les prolongements possèdent des poids différents,
l'expression pondérée conserve cette différence structurelle. La
constante de Boltzmann peut être définie ici par sa fonction : c'est le
coefficient dimensionnel qui transporte l'information de la
distribution admise vers la droite d'entropie et relie, sous une
température déclarée, cette information à l'énergie.

L'action intervient avant la publication thermique. La section
\(\hbar_a\) convertit la fréquence de l'évolution en énergie ;
la section \(h_a\) exprime l'action accumulée pendant un tour.
L'égalité \(k_B\Theta_{{\rm clk},a}\log3=h_a/T_{108}\)
réunit les deux fonctions : énergie de la période et information locale de
la dynamique tritique. Une identification thermodynamique conserve
en outre l'état statistique, la température et l'opération physique
d'échange ou d'effacement étudiée.
